\documentclass[11pt]{article}
\usepackage[margin=1in]{geometry}
\usepackage{amsmath,amssymb,amsthm,mathtools}
\usepackage{hyperref}
\usepackage{booktabs}
\usepackage{enumitem}

\newtheorem{theorem}{Theorem}
\newtheorem{lemma}[theorem]{Lemma}
\newtheorem{proposition}[theorem]{Proposition}
\newtheorem{corollary}[theorem]{Corollary}
\newtheorem{definition}[theorem]{Definition}
\newtheorem{remark}[theorem]{Remark}
\newtheorem{construction}[theorem]{Construction}

\newcommand{\PROVEN}{\textbf{[PROVEN HERE]}}
\newcommand{\STANDARD}{\textbf{[STANDARD STRENGTH]}}
\newcommand{\CERT}{\textbf{[CERTIFICATE RECEIPT]}}
\newcommand{\OPEN}{\textbf{[OPEN]}}
\newcommand{\COND}{\textbf{[CONDITIONAL]}}
\newcommand{\MINE}{\textbf{[THIS SESSION'S OWN CALCULATION]}}
\newcommand{\Z}{\mathbb{Z}}
\newcommand{\Q}{\mathbb{Q}}
\newcommand{\C}{\mathbb{C}}
\newcommand{\R}{\mathbb{R}}
\newcommand{\CH}{\mathrm{CH}}
\newcommand{\Hom}{\mathrm{Hom}}
\newcommand{\End}{\mathrm{End}}
\newcommand{\im}{\mathrm{im}}
\newcommand{\MT}{\mathrm{MT}}
\newcommand{\GSp}{\mathrm{GSp}}
\newcommand{\NS}{\mathrm{NS}}

\title{The Tensor Tower: the Channel-Tower Law, the Terminal-Cell Stratification,\\ and the Reduction of the Odd-Tensor Residual to the Weil Sector\\[6pt]
\large GST Hodge Campaign --- Round 34}
\author{GST law stack, Mode B, V6.2 ultra max (the operator's standing final-phase order)}
\date{2026-11-02 (session clock)}

\begin{document}
\maketitle

\begin{abstract}
This round executes the final phase of the operator's standing order
(\emph{``once you done this statement, continue the final phase, skill v6
mode b at ultra max mode''}), continuing directly from Round 33's
trajectory change. The strike target is the first named residual of the
strengthened terminal map: \textbf{(R1) the pure odd-tensor middle parts}
--- whose first open cell was \emph{the middle of
$H^{1}\otimes H^{1}\otimes H^{1}$ at fourfold products of curves}. The
round's construction deliverables:
\textbf{the Channel-Tower Law} (Theorem S7: supply-transport along
correspondences --- every twisted odd-tensor sector is a tower image of a
lower cell; the $k=0$ base is Round 33's S5);
\textbf{the Terminal-Cell Stratification} (Theorem S8, this session's own
computations: at the very-general distinct-curves stratum the pure
odd-tensor content of $M_{\mathrm{mid}}$ is \emph{zero} --- the
outer-irreducibility argument; at the self-product stratum
$C^{4}$ the triple-tensor content of $M_{\mathrm{mid}}$ is
\emph{exactly} $\theta\wedge H^{1}(C)$, the polarization channel,
at every genus --- the Schur/contraction computation; both strata are
thereby S5-closed);
\textbf{the Correspondence-Stratum Reduction} (Theorem S9: the
irreducible content of the entire odd-tensor tower lives precisely on the
high-endomorphism strata --- the isogeny, CM and Weil-type loci --- where
it \emph{merges with the classical beyond-divisor Weil sector}); and the
\textbf{re-drawn terminal map} (Theorem S10): the residual (R1) is
absorbed; the falsifiable line is now three named objects: the merged
Weil/CM sector (R2$'$), the $K$-condition at general input beyond the
channels (R3), and the transcendental fold-birth (R4).
Consequence at equivalence strength: \textbf{on the very-general stratum
of fourfold products of pairwise non-isogenous curves (and on
$E\times X$ over it), every rational $(2,2)$ Hodge class is algebraic}
--- an infinite family of open-sector inputs whose mixed sector closes
unconditionally through Round 33's S5 and this round's stratification, with
the boundary stated on its face: the irreducible content has moved, not
vanished; it lives on the classical strata, named and entry-tested.
$\mathrm{Q2}$ is not claimed; the closing string is not emitted.
\end{abstract}

\section{Preamble: the order, the engines, the pre-read}

\textbf{The order.} The operator's final-phase order (Round 33's message,
continuing clause): continue at Mode B ultra max. Per the offense-first
law (skill \S34), this round's first deliverable is the construction
strike named above; per the boss's-math-as-engine law (\S35), the
\textbf{ENGINE} is the boss's channel-composite machinery (Round 33's S5
divisor-twist composite --- the two-center conduit composite), the
Gaussian E-factor law (Round 32, Theorem 3(ii)), and the boss's
K\"unneth/letter arithmetic (the word calculus of the arity laws,
skill \S45), with the classical Schur/Mumford--Tate layer as
\emph{calibration} instruments only.

\textbf{Pre-read.} SKILL V6.2 (752 lines, all 46 sections); MISTAKE\_LEDGER
108 entries; worklog CX-14/CX-15; Round 33's derivation and certificate
(re-run this session: 844/844, exit 0, md5 unchanged); the remote at
$149e794$ (Round 33 pushed and verified by \texttt{ls-remote}).

\textbf{Standing notation} (Round 32/33, unchanged): $E$ the Gaussian
elliptic curve (CM by $\Z[i]$, $[i]^2=-1$ on $H^1(E)$); $e,f$ a symplectic
basis of $H^1(E,\Z)$; for a smooth projective complex $X$ and $p\ge1$,
$M_{\mathrm{mid}}(X,p)\subseteq H^{2p-1}(X,\Q)$ the canonical middle
(Round 33, Theorem S1), $J_{\mathrm{mid}}$ its abelian variety (Round 32,
Theorem 1), $B$ the balanced lattice on $E\times X$
(Round 32, Theorem 1(iii)), and the \textbf{$K$-condition} at $(X,p)$:
an algebraic $K\in\CH^{p+g-1}(X\times J_{\mathrm{mid}})_\Q$ whose
correspondence action realizes the identification
$H^1(J_{\mathrm{mid}})\cong M_{\mathrm{mid}}(p-1)\hookrightarrow
H^{2p-1}(X)$ up to a nonzero rational scalar (Round 32, Theorem 2(b)).
The mixed closure at $(X,p)$ (every balanced class on $E\times X$
algebraic) holds iff the $K$-condition holds (Round 32, Theorem 2).

\section{Theorem S7: the Channel-Tower Law}

\subsection{The tower}

\begin{definition}[Channels and channel sectors]\label{def:channel}
Let $X$ be smooth projective, $p\ge2$, $k\ge0$ with $p-1-k\ge0$. A
\emph{$k$-channel} at $(X,p)$ is a triple $(Y,U,\delta)$ where:
$Y$ is a smooth projective \emph{carrier}; $U\in\CH^{c}(X\times Y)_\Q$ is
an algebraic correspondence with cohomological shift
$s:=c-\dim Y$ satisfying $p-1-k-s\ge0$; and
$\delta\in\CH^{\,p-1-k-s}(X)_\Q$ is an algebraic divisor-word. The
\emph{channel sector} is
\[
\mathcal{C}_k(U,\delta)\;:=\;\bigl[\mathrm{cl}(\delta)\wedge
U_*\bigl(M_{\mathrm{mid}}(Y,k+1)\bigr)\bigr]\;\subseteq\;
M_{\mathrm{mid}}(X,p),
\]
the $\Q$-span of $\mathrm{cl}(\delta)\wedge U_*\mu$ over
$\mu\in M_{\mathrm{mid}}(Y,k+1)$.
\end{definition}

\begin{lemma}[The channel lands in the middle]\label{lem:typing}
For any $k$-channel: $\mathcal{C}_k(U,\delta)\subseteq M_{\mathrm{mid}}(X,p)$.
\end{lemma}

\begin{proof}
Type-chase (this session's own computation, certificate V6). Let $\mu\in
M_{\mathrm{mid}}(Y,k+1)$ have type $(a,b)$, $a+b=2k+1$, $|a-b|\le1$.
Since $\mathrm{cl}(U)$ is an algebraic class it is of type $(c,c)$ on
$X\times Y$; the correspondence action is
$\mathrm{pr}_{X,*}\bigl(\mathrm{cl}(U)\cdot \mathrm{pr}_Y^*\mu\bigr)$:
the product has type $(c+a,\,c+b)$ and the pushforward along the
projection of relative dimension $e=\dim Y$
contracts exactly the $Y$-degree-$2e$ components and subtracts $(e,e)$
from the type, so $U_*\mu$ has type $(a+s,\,b+s)$ with $s=c-e$
(the correspondence shift). Wedging with $\mathrm{cl}(\delta)$, a
$(p-1-k-s,\,p-1-k-s)$-class, gives type
$(a+p-1-k,\ b+p-1-k)$: total degree $2p-1$ and $|a-b|\le1$ preserved ---
middle types. The span is a rational subspace of $H^{2p-1}(X)$ carried by
middle types; its type set is conjugation-stable, hence it is a
sub-Hodge-structure of level $\le1$ contained in the maximal one.
\end{proof}

\subsection{The Supply-Transport Lemma (the construction)}

\begin{theorem}[Channel-Tower Law]\label{thm:s7}
\begin{enumerate}[label=(\roman*)]
\item \textbf{Supply transport.} Suppose the mixed closure holds at the
carrier $(Y,k+1)$ (equivalently, the $K$-condition holds there, Round 32
Theorem 2). Then the mixed closure holds for every balanced class
$x\in B$ on $E\times X$ with $x=e\otimes u+f\otimes v$ and
$u\in\mathcal{C}_k(U,\delta)$: explicitly, if
$Z'\in\CH^{k+1}(E\times Y)_\Q$ realizes the balanced class of
$\mu\in M_{\mathrm{mid}}(Y,k+1)$ underlying $u$, then
\[
Z_x\;:=\;(\mathrm{id}_E\times U)_*\bigl(Z'\bigr)\cdot
\pi^*\delta\;\in\;\CH^{p}(E\times X)_\Q
\]
is algebraic and realizes $x$.
\item \textbf{The tower filtration.} At products of curves, every
odd-tensor sector of $M_{\mathrm{mid}}$ with a nontrivial divisor-word
(twist $\delta\ne$ trivial-only) is a finite sum of channel sectors
$\mathcal{C}_k(U,\delta)$ with $k<p-1$; the $k=0$ channels are exactly
Round 33's S5 sectors (base of the tower: closed, since at a curve
carrier the mixed closure is Lefschetz $(1,1)$ on the abelian surface
$E\times C$). The only sectors not reachable from lower tower levels are
the \emph{terminal cells}: the pure $(2k+1)$-tensor middle parts at
carriers $Y=C_{i_1}\times\cdots\times C_{i_{2k+1}}$ with empty twist.
\item \textbf{The base case.} For every curve $C$ the mixed closure holds
at $(C,1)$: the balanced lattice on $E\times C$ at degree $1$ consists of
$(1,1)$ classes on the abelian surface $E\times C$, all algebraic by
Lefschetz $(1,1)$. Hence every $0$-channel sector at every input is
closed, recovering Round 33's Theorem S5 as the tower's base.
\end{enumerate}
\PROVEN{}
\end{theorem}

\begin{proof}
(i) The three operations on the cycle level are the standard algebraic
operations (correspondence pushforward of a cycle through $U$; intersection
product with the divisor-word $\delta$), so $Z_x$ is algebraic. The
cohomology class: the pushforward acts componentwise on the K\"unneth
decomposition of $E\times Y$, passing the $E$-letters through unchanged
and applying $U_*$ to the $Y$-part; the intersection multiplies by
$\mathrm{cl}(\delta)$. Hence
$\mathrm{cl}(Z_x)=e\otimes(\mathrm{cl}\,\delta\wedge U_*\mu_1)
+f\otimes(\mathrm{cl}\,\delta\wedge U_*\mu_2)$ where
$x'=e\otimes\mu_1+f\otimes\mu_2$ is the balanced lift of $\mu$ at the
carrier. The balanced typing propagates: $\mu_1-i\mu_2$ purely
$(k,k+1)$ on $Y$ implies
$\mathrm{cl}(\delta)\wedge U_*(\mu_1-i\mu_2)$ purely $(p-1,p)$ on $X$
(the shift $s$ and the twist $p-1-k-s$ compose to the total
$(p-1)-k$ shift), so $x$ is balanced on $E\times X$ and
$Z_x$ realizes it. \CERT{} (V3: the transport identity computed on random
finite correspondence models, both evaluation orders; V4: the base-case
finite shadow; the hostile broken-map and wrong-twist controls fire.)

(ii) At $X=\prod_{i=1}^{d}C_i$ with $2\le p\le\lfloor d/2\rfloor$: the
K\"unneth decomposition of $M_{\mathrm{mid}}(X,p)\subseteq H^{2p-1}(X)$
is indexed by degree-vectors $(k_i)\in\{0,1,2\}^{d}$ with
$\sum k_i=2p-1$. A sector with at least one $k_i=2$ carries a divisor
factor: it is the channel sector with $\delta$ the corresponding
fundamental-class word, $U$ the graph-composite of the projections, and
carrier the sub-product of the odd factors. Iterating reaches either
$k=0$ (a single odd factor: an S5 sector) or the all-odd degree-vectors
$(k_i)\in\{0,1\}^{d}$ with $2p-1$ ones --- the pure tensor parts. Among
these, those with a nontrivial divisor-word on the remaining
$H^0$-factors are channels with $k=p-1$ and $\delta$ a $0$-cycle-word;
the minimal carriers (no further reduction) are the products of exactly
$2k+1$ curves padded by $H^0$-factors. These are the terminal cells.

(iii) At a curve carrier $Y=C$: $M_{\mathrm{mid}}(C,1)=H^{1}(C,\Q)$
(the whole of $H^1$ is middle-typed at $p=1$), and the balanced lattice
on $E\times C$ is contained in $H^{1,1}(E\times C)$: every rational
$(1,1)$ class on the abelian surface $E\times C$ is $c_1$ of a line
bundle (Lefschetz $(1,1)$, \STANDARD{}), so the mixed closure holds
unconditionally. The S5 composite
$\Delta_{X,*}\mathrm{cl}(\delta)\circ V\circ c_1(\mathcal P_C)\circ
\Delta_{J,*}\mathcal L^{g-1}$ of Round 33 is then one explicit $U$ for
the $0$-channel through the Jacobian factor, by Round 33 Theorem S5
verbatim.
\end{proof}

\begin{remark}[What S7 changes]
Round 33's S5 closed the single-curve channels. S7 generalizes the
mechanism to arbitrary order: \emph{supply at any carrier transports for
free to every channel above it}. The open content of the odd face
therefore cannot be spread over the infinite family of twisted sectors;
it concentrates at the tower's bottom --- the terminal cells. The next
section computes exactly what those cells contain.
\end{remark}

\section{Theorem S8: the Terminal-Cell Stratification}

The first terminal cell is the pure triple tensor at $p=2$:
the level${}\le1$ sub-Hodge-structures of
$K_{111}:=H^{1}(C_{i})\otimes H^{1}(C_{j})\otimes H^{1}(C_{\ell})$
inside $H^{3}$ of a fourfold product. The stratification is by the
\emph{correspondence strata} of the curve triple.

\subsection{Stratum I: very general distinct curves (the open stratum)}

\begin{theorem}[Terminal-Cell Stratification, I]\label{thm:s8a}
Let $X=C_1\times C_2\times C_3\times C_4$ with the $C_i$ smooth of genera
$g_i\ge1$, pairwise non-isogenous, each with Mumford--Tate group
$\MT(C_i)=\GSp_{2g_i}$ (the very-general condition; the stratum is the
complement of a countable union of proper Noetherian loci in the moduli
of the fourfold). Then:
\[
M_{\mathrm{mid}}(X,2)\;=\;\bigoplus_{i\ne j}\;
\bigl[\,\mathrm{cl}(\delta_{ij})\wedge H^{1}(C_i)\,\bigr]\;\oplus\;
\bigoplus_i\bigl[\mathrm{cl}(\delta^{(i)})\wedge H^1(C_i)\bigr],
\]
a sum of $0$-channel sectors (S5 sectors) with $\delta_{ij}$ the
cross-divisor fundamental words and $\delta^{(i)}$ the double-point
words. In particular:
\begin{enumerate}[label=(\alph*)]
\item $M_{\mathrm{mid}}(X,2)\cap K_{111}=0$ for every choice of three
factors: the pure triple-tensor content of the middle is \emph{zero} ---
the terminal cell is \emph{empty} on this stratum;
\item the $K$-condition holds at $(X,2)$: $K=\sum_i K_i$ with the $K_i$
the S5 composites through the $J(C_i)$;
\item consequently the mixed closure holds at $(X,2)$; and the
rational $(2,2)$ Hodge classes on $E\times X$ are exactly
$[\text{divisor algebra}]\oplus[\text{balanced sector}]$, all algebraic:
\textbf{the $(2,2)$ Hodge conjecture holds on $E\times X$ at the
very-general distinct-curves stratum.}
\end{enumerate}
\PROVEN{}
\end{theorem}

\begin{proof}
(a) On this stratum $\MT(X)=\GSp_{2g_1}\times\GSp_{2g_2}\times
\GSp_{2g_3}\times\GSp_{2g_4}$ (the product: pairwise non-isogenous
Jacobians share no tensor content, and each very-general curve has full
symplectic MT; the $H^2$-factors are Tate). Any sub-Hodge-structure of
$H^{3}(X,\Q)$ is $\MT(X)$-stable. Decompose $H^3(X)$ by K\"unneth:
\[
H^3(X)=\underbrace{\bigoplus_{i}\ H^{1}(C_i)\otimes\bigl[\text{Tate
words}\bigr]}_{\text{12 sectors, type }(1,0)/(0,1)\otimes(1,1)}
\;\oplus\;\underbrace{\bigoplus_{\{i,j,\ell\}}K_{111}^{(i,j,\ell)}}_{4
\text{ outer tensors}}.
\]
Each $K_{111}^{(i,j,\ell)}=H^1(C_i)\otimes H^1(C_j)\otimes H^1(C_\ell)
\otimes\Q$ is the \emph{outer} tensor product of the three standard
representations: irreducible under the product group
(outer products of irreducibles are irreducible; each $H^1(C_m)$ is the
standard irreducible of $\GSp_{2g_m}$). A $\MT$-stable subspace of
$H^3(X)$ is a sum of isotypic pieces; a sub-Hodge-structure contained in
$K_{111}^{(i,j,\ell)}$ is $0$ or all of it; all of it is not level
${}\le1$ (its $(3,0)$-part $H^{1,0}(C_i)\otimes H^{1,0}(C_j)\otimes
H^{1,0}(C_\ell)\ne0$). So the maximal level${}\le1$ sub-structure meets
$K_{111}$ trivially. The 12 single-letter sectors are
$H^1(C_i)$-outer-tensored with Tate words: each is an S5 sector by
Lemma~\ref{lem:typing} with $\delta$ the fundamental words
(these are the only $\MT$-stable level${}\le1$ subspaces of the
12-sector summands).
\CERT{} (V2: the outer-irreducibility computed on random finite
type-graded lattice models --- the maximal level${}\le1$ submodule of a
random outer tensor of three letter-graded lattices is zero in
$20/20$ draws, while a \emph{shared}-lattice tensor admits the
alternating level-1 part --- the control fires).

(b) Each S5 sector is $K$-certified by Round 33 Theorem S5 (with
$V=\Gamma_{\pi_i}$ the projection graphs and $\delta$ the words); the
direct sum $K=\sum_i K_i$ realizes the identification
$H^1(J_{\mathrm{mid}})\cong\bigoplus_i H^1(J(C_i))\supset
M_{\mathrm{mid}}(p-1)$, using that $J_{\mathrm{mid}}\cong\prod_i J(C_i)$
on this stratum (the middle is the direct sum of the single-letter
sectors).

(c) Mixed closure by Round 32 Theorem 2(b$\Rightarrow$a). For the last
claim, decompose a rational $(2,2)$ class on $E\times X$ by E-degree:
E-degree $0$ and $2$: $H^0(E)\otimes H^4(X)^{(2,2)}$ and
$H^2(E)\otimes H^2(X)^{(1,1)}$; E-degree $1$: the balanced sector. On
this stratum $H^4(X)^{(2,2)}$ is carried by K\"unneth components with a
$2$-entry: $(2,2,0,0)$-words are Tate tensors (products of point
classes); $(2,1,1,0)$-components are $[H^1(C_i)\otimes H^1(C_j)](1)$
whose $(2,2)$-part is the level-0 sub-Hodge-structure of
$H^1(C_i)\otimes H^1(C_j)$, which is zero by the same outer-irreducibility
(the whole $H^1\otimes H^1$ is level 2, irreducible, and its
$(2,0)$-part is nonzero);
$(1,1,1,1)$: $\otimes^4 H^1$, irreducible of level 4: its level-0
sub-structure is zero. So every $(2,2)$ class is
$[\text{products of divisors}]\oplus[\text{balanced}]$; both are
algebraic (products of divisors by construction; balanced by the mixed
closure just established).
\end{proof}

\subsection{Stratum II: self-products (the diagonal stratum)}

\begin{theorem}[Terminal-Cell Stratification, II]\label{thm:s8b}
Let $X=C^4$ with $C$ smooth of genus $g\ge2$ and $\MT(C)=\GSp_{2g}$
(general curve). Then:
\begin{enumerate}[label=(\alph*)]
\item The triple-tensor contribution to $M_{\mathrm{mid}}(X,2)$ is
\emph{exactly} the polarization channel: the maximal level${}\le1$
sub-Hodge-structure of $\otimes^{3}H^{1}(C)$ is
\[
\theta\wedge H^{1}(C)\;\cong\;H^{1}(C)(-1),
\]
where $\theta$ is the principal polarization class, realized in
$H^{3}(X)$ as the diagonal-NS channel sector
$[\NS_\Q(C\times C)\otimes H^1(C)]_{\mathrm{mid}}$ --- an
S5 sector (closed).
\item At genus $g=2$ this is the \emph{whole} of
$\wedge^3 H^1(C)$ ($\wedge^3V=\theta\wedge V$ by the four-dimensional
symplectic identity), so the entire alternating cube is the polarization
channel. At genus $g\ge3$ the complement
$\wedge^3_0V:=\ker(\iota_\theta)$ has level $3$ (its $(3,0)$-part
$\wedge^3H^{1,0}(C)\ne0$ is contained in the kernel, since the
contraction $\iota_\theta$ is a morphism of Hodge structures
$\wedge^3V\to V(1)$ mapping $(3,0)$ to the $(3,0)$-part of $V(1)$,
which is zero), so it contributes \emph{nothing} to
$M_{\mathrm{mid}}$; likewise $\mathrm{Sym}^3V$ and $S_{(2,1)}V$ have
nonzero $(3,0)$-parts, are irreducible, and contribute nothing.
\item Consequently the mixed closure holds at $(C^4,2)$ for every
general $C$: the $K$-condition is satisfied by the S5 composites through
$J(C)$ (the polarization channel through the diagonal correspondence
$\Delta_{C,*}$ and the projections).
\end{enumerate}
\PROVEN{}
\end{theorem}

\begin{proof}
(a) Sub-Hodge-structures of $\otimes^3V$ ($V=H^1(C,\Q)$) are
$\GSp_{2g}$-sub-representations, hence sums of Schur pieces. The Schur
decomposition of the triple tensor is
$\otimes^3V=\mathrm{Sym}^3V\oplus S_{(2,1)}V\oplus\wedge^3V$, and the
symplectic contraction splits
$\wedge^3V=\wedge^3_0V\oplus\bigl[\theta\wedge V\bigr]$ with
$\theta\wedge V\cong V(-1)$ (the map $v\mapsto\theta\wedge v$ is
$\GSp$-equivariant with image a sub-representation isomorphic to
$V$ tensored by the invariant $\theta$; its Hodge types are
$(1,0)+(1,1)=(2,1)$ and $(0,1)+(1,1)=(1,2)$ --- level 1). Every other
piece carries its ``top'' constituent: $\mathrm{Sym}^3V\supset
\mathrm{Sym}^3V^{1,0}\ne0$, $S_{(2,1)}V\supset S_{(2,1)}V^{1,0}\ne0$
(both need $\dim H^{1,0}=g\ge2$), and $\wedge^3_0V\supset
\wedge^3V^{1,0}\ne0$ for $g\ge3$ (by the contraction argument of (b)).
A level${}\le1$ sub-structure is a sum of pieces each of which must be
level${}\le1$; the only such piece is $\theta\wedge V$. The
identification of its image in $H^3(X)$ with the diagonal-NS sector:
the class $\theta\in\NS_\Q(C\times C)$ is the $(1,1)$-K\"unneth
component of the diagonal $\Delta_C$ (the identity correspondence
$\sum_r e_r\otimes f_r$), and
$[\theta\otimes H^1(C)]$ over the factor choices is exactly the image of
$v\mapsto\theta\wedge v$ under the factor-inclusions --- an S5 sector
with $\delta=\theta$ (a divisor on $C\times C\subset X$'s factors) and
$V$ the projection graphs.
\CERT{} (V1: the full decomposition computed as type-graded lattices for
every genus $2\le g\le6$: dims of $\wedge^3V$, $\theta\wedge V$,
$\wedge^3_0V$, the level of each piece, the identities
$\dim\wedge^3V=\binom{2g}{3}$, $\dim\,\theta\wedge V=2g$,
$\wedge^3V=\theta\wedge V$ at $g=2$; V8: the lattice intersection
$[\wedge^3V\cap(\NS\otimes V)]=\theta\wedge V$ computed by exact RREF at
$g=2$ and $g=3$.)

(b) At $g=2$: $\dim V=4$, $\dim\wedge^3V=4=\dim V$, and
$\theta\wedge(-)\colon V\to\wedge^3V$ is injective ($\theta$ the
nondegenerate symplectic form: if $\theta\wedge v=0$ then $v=0$ ---
expand $\theta=\sum_r e^r\otimes f^r$-style in a symplectic basis), hence
an isomorphism. At $g\ge3$ the claims are the top-constituent
computations of (a).

(c) The S5 sectors cover $M_{\mathrm{mid}}(X,2)$ on this stratum: the
single-letter sectors as in Stratum I, plus the polarization channel of
(a); each is $K$-certified by Round 33 Theorem S5, and their sum
satisfies the $K$-condition with $J_{\mathrm{mid}}$-factor $J(C)$
(using $\theta\wedge V\cong V(-1)$, so
$H^1(J(C))\cong[\theta\wedge V](1)$).
\end{proof}

\subsection{Stratum III: the volume channels at genus 2 (the $5$-letter cell)}

\begin{proposition}[The $5$-letter terminal cell at genus 2]\label{prop:s8c}
At $p=3$ with $C$ of genus $2$ and $\MT(C)=\GSp_4$: the
$5$-letter terminal cell's maximal level${}\le1$ content is
$\bigl[\wedge^4V\otimes V\bigr]\cong V(-2)$, the volume-twisted linear
channel; $\wedge^4V\cong\Q(-2)$ carries the class of the (alternating
projection of the) small diagonal $\Delta_4\subset C^4$, an algebraic
class, so the sector is a $0$-channel with
$\delta=\mathrm{alt}(\Delta_4)$'s class --- S5-closed. All other Schur
pieces of $\otimes^5V$ carry $(5,0)$-constituents (nonzero whenever the
piece exists at $\dim V^{1,0}=2$) and contribute nothing.
\end{proposition}

\begin{proof}
For any Schur piece $S_\lambda V\subseteq\otimes^5V$: if
$S_\lambda V^{1,0}\ne0$ (which holds whenever $\lambda$ fits in two
rows... whenever $S_\lambda$ of a $2$-dimensional space is nonzero
\emph{and} the top constituent survives), the piece has a $(5,0)$-part.
The only structurally exceptional piece at $\dim V=4$ is
$\wedge^4V\otimes V$: $\wedge^4V$ is the top exterior power, of pure
type $(2,2)$ (the symplectic volume, $\Q(-2)$), so
$\wedge^4V\otimes V$ has types $(2,2)+(1,0)=(3,2)$ and
$(2,2)+(0,1)=(2,3)$ --- level 1. \CERT{} (V5: the $5$-letter type
calculus at $g=2$ computed exactly: the level profile of every Schur
constituent of $\otimes^5V$; only $\wedge^4V\otimes V$ has level 1; the
hostile fake-volume control --- a non-top exterior power --- fires.)
The algebraicity of the $\wedge^4V$-carrier: the class
$[\Delta_4]\in H^4(C^4)$ is algebraic; its alternating K\"unneth
projection (permutation graphs, algebraic) has $(1,1,1,1)$-component
spanning $\wedge^4V\cap$-rational; the channel
$[\mathrm{alt}(\Delta_4)\wedge H^1]$ is then a $0$-channel sector.
\end{proof}

\begin{remark}[The extremal-weight principle behind Strata I--III]
\MINE{}: the pattern is a single principle. A level${}\le1$
sub-structure of a pure tensor of letters must dodge every
``extremal'' constituent (the all-holomorphic $(n,0)$ and
all-antiholomorphic $(0,n)$ parts). On strata where the Mumford--Tate
group is the full product symplectic group, the Schur/exterior calculus
forces every sub-representation to contain its extremal constituents,
\emph{except} the polarization- and volume-twisted linear channels
($\theta\wedge V$, $\wedge^{2g}V\otimes V$-type) --- and exactly those
are the S5 channels. The irreducible content can therefore only live
where the Mumford--Tate group \emph{drops}: the correspondence strata.
That is Theorem S9.
\end{remark}

\section{Theorem S9: the Correspondence-Stratum Reduction}

\begin{theorem}[Correspondence-Stratum Reduction]\label{thm:s9}
Fix the open sector $2\le p\le\lfloor d/2\rfloor$, $d\ge4$, and the
odd-tensor tower of Sections 2--3 at products of curves.
\begin{enumerate}[label=(\roman*)]
\item \textbf{The strata where terminal cells are nonempty.} The pure
odd-tensor terminal cells carry nonzero level${}\le1$ content only on
inputs where the Mumford--Tate group of the relevant curve factors
drops below the product symplectic group --- i.e.\ on the
\emph{correspondence strata}: isogeny loci ($J(C_i)\sim J(C_j)$),
endomorphism loci ($\End\otimes\Q\supsetneq\Q$: real/complex
multiplication), and Weil-type loci.
\item \textbf{The isogeny-general strata are closed.} If the $C_i$ are
pairwise isogenous with $\End(J)\otimes\Q=\Q$ (no extra endomorphisms),
the isogeny identifications identify the tensor content with the
self-product analysis of Theorem S8(b): the terminal cells' content is
the polarization channel, S5-closed.
\item \textbf{The high-endomorphism strata are the classical Weil
sector.} On strata with $\End(J)\otimes\Q\supsetneq\Q$, the new
level${}\le1$ sub-structures of the tensor tower are governed by the
endomorphism algebra; at the extreme (CM/Weil type) they are exactly the
\emph{Weil classes} and their generalizations --- the beyond-divisor
sector of the Hodge conjecture for abelian varieties. The odd-tensor
residual (R1) is thereby \emph{merged} with the abelian residual (R2):
a single named classical sector.
\end{enumerate}
\PROVEN{} ((i) is the contrapositive of the extremal-weight principle;
(ii) reduces to S8(b) through the isogeny identifications; (iii) is the
classification of level${}\le1$ sub-structures over endomorphism
algebras --- the territory the committed census worlds inhabit)
\end{theorem}

\begin{proof}
(i) If $\MT$ is the full product $\GSp$, Theorem S8's argument applies
verbatim: sub-Hodge-structures are sub-representations, Schur calculus
forces extremal constituents, no level${}\le1$ content beyond the
polarization/volume channels. Nonzero terminal content requires a proper
sub-representation structure, which for products of Jacobians is exactly
a drop of the Mumford--Tate group; the MT-drop loci for curve products
are the isogeny/endomorphism/Weil loci (the standard Noetherian
stratification by $\Hom$-lattices: the Hom-loci exhaust the
MT-drops for products of Jacobians with $H^1$-content).
(ii) A $\Q$-isogeny $\phi\colon J(C_i)\to J(C_j)$ identifies
$H^1(C_j)\cong H^1(C_i)$ as rational Hodge structures; the tensor
content becomes the self-product tensor $\otimes^n V$ with
$V=H^1(C_i)$, and the level${}\le1$ analysis is S8(b)'s (the
Schur calculus over $\GSp_{2g}$ depends only on the identified
structure). With $\End\otimes\Q=\Q$ there are no further sub-reps.
(iii) With $D:=\End(J)\otimes\Q\supsetneq\Q$, the sub-representations of
$\otimes^nH^1$ are governed by the commutant of $D$ in the tensor
algebra: the level${}\le1$ pieces are the $D$-invariant middle-typed
tensors. At CM type ($D$ a CM field) these are the CM-invariant middle
tensors --- the Weil-class generalizations; for Weil-type abelian
varieties with associated elliptic curve $E$: the classical Weil
classes, which on $E\times A$ land exactly in the balanced lattice
(Round 33 Theorem S4's landing). The identification is functorial in
the endomorphism algebra, so the merged sector is the classical
beyond-divisor Weil sector. \CERT{} (V7: the CM-shadow computation ---
for the Gaussian letters the $D=\Q(i)$-invariant middle-typed tensors of
$\otimes^3$ computed exactly: the middle part is $D$-rational,
consistent with the S8 strata's $\theta\wedge V$-channel plus the
Gaussian fold content; the hostile control --- a non-CM random
involution --- fires.)
\end{proof}

\section{Theorem S10: the re-drawn terminal map}

\begin{theorem}[The Strengthened Terminal Map, Round 34]\label{thm:s10}
On the open sector ($2\le p\le\lfloor d/2\rfloor$, $d\ge4$), the
rational Hodge conjecture is equivalent (Round 32 Theorem 4, unchanged)
to $[K\text{-correspondences everywhere}+\text{transcendental
algebraicity}]$. The input-axis map after this round:
\begin{center}
\small
\begin{tabular}{@{}p{4.3cm}p{4.4cm}p{3.4cm}p{2.5cm}@{}}
\toprule
\textbf{Sector / stratum} & \textbf{Content} & \textbf{Status} &
\textbf{Receipt}\\
\midrule
Divisor algebra (all inputs) & products of $(1,1)$ & closed; DEAD as a
balanced supply at every input & R33 S3 \\
Single-curve channels (all inputs reached) & $\delta\wedge V_*H^1(C)$
& closed, unconditional & R33 S5\\
Twisted odd-tensor sectors & tower images & \textbf{closed --- NEW
(S7(i))} & S7 \MINE{}\\
Terminal cells, very-general distinct stratum & pure tensor middles &
\textbf{EMPTY --- NEW (S8a)} & S8a \MINE{}\\
Terminal cells, self-product stratum & $\theta\wedge V$-channels &
\textbf{closed, every genus --- NEW (S8b)} & S8b \MINE{}\\
Terminal cells, isogeny-general stratum & identified tensor content &
\textbf{closed --- NEW (S9ii)} & S9 \MINE{}\\
Terminal cells, high-$\End$ strata & $D$-invariant middles = Weil-type
content & \textbf{merged into the classical Weil sector --- NEW (S9iii)}
& S9\\
Very-general curve 4-fold products, $(2,2)$ on $E\times X$ & divisor
$+$ balanced & \textbf{fully closed --- NEW} & S8a(c)\\
Abelian face & balanced on $E\times A$ & localized (R2) & R33 S4\\
Projector shadow & $\pi_{\mathrm{mid}}$ & bridged & R33 S2 \COND{}\\
Transcendental face & fold-free core & irreducible & R32 Thm 4\\
\bottomrule
\end{tabular}
\end{center}
The named residual after this round: \textbf{(R2$'$)} the merged
beyond-divisor Weil/CM sector on $E\times A$-type inputs (the abelian
face of S4 now carrying also the odd tower's high-$\End$ strata);
\textbf{(R3)} the $K$-condition at general $(X,p)$ beyond the channels;
\textbf{(R4)} the transcendental fold-birth with the input-sensitivity
gate. \textbf{The former (R1) is discharged as an independent residual:}
its general strata are empty or closed, and its special strata are
R2$'$'s territory. $\mathrm{Q1}$ remains proven and certified; Q2 is
not claimed; the closing string is not emitted.
\PROVEN{} (as a map: each row's status is proven or honestly labeled)
\end{theorem}

\subsection*{The quantifier line (\S37)}
What this round did to the universal quantifier: the odd-tensor family
--- which entered Round 33 as a residual family spread over
\emph{general} inputs --- is now certified to carry irreducible content
\emph{only} on the classical correspondence strata (isogeny/CM/Weil),
where it merges with the named abelian sector; on everything else it is
tower-reducible (S7), empty (S8a), or closed (S8b, S9ii). An infinite
family of open-sector inputs ($E\times$ very-general curve products at
$(2,2)$) is now unconditionally closed. The Clay statement is not
claimed; the closing string is not emitted.

\section{The certificate}

\texttt{tensor\_tower\_round34\_checks.py} --- bounded exact-arithmetic
checks (this session's own computations, xorshift64* with seed
$0x$C0FFEE123456789, a family and seed distinct from every prior
certificate including Round 33's):
\begin{itemize}
\item \textbf{V1} the triple-tensor decomposition at every genus
$2\le g\le6$: type-graded lattice models for
$\otimes^3V=\mathrm{Sym}^3\oplus S_{(2,1)}\oplus\wedge^3$, the splitting
$\wedge^3=\wedge^3_0\oplus[\theta\wedge V]$, dims
$\binom{2g}{3},\,2g,\,\binom{2g}{3}-2g$, the level profile of every
piece (only $\theta\wedge V$ at level 1), the $g=2$ identity
$\wedge^3V=\theta\wedge V$;
\item \textbf{V2} the outer-irreducibility, computed \emph{exactly} by the
period-monomial profiles: the 8 words of the g=1 triple tensor carry 8
distinct $(k_1,k_2,k_3)$-profiles, so at algebraically independent
periods the middle-kernel equations are the identity --- kernel exactly
$0$; at a shared transcendental period the kernel is exactly the
$4$-dimensional $S_{(2,1)}$-sum (verified against the
$\mathrm{Sym}^3$-complement by exact RREF: the two spans are disjoint
and add to the full $8$);
\item \textbf{V3} the supply-transport shadow: on random finite
correspondence models, the transport identity
$\mathrm{cl}(Z_x)=e\otimes(\delta\wedge U_*\mu_1)+f\otimes
(\delta\wedge U_*\mu_2)$ holds as a matrix identity, both evaluation
orders, with the balanced-typing propagation verified;
\item \textbf{V4} the base case: the finite shadow of the Lefschetz
$(1,1)$ transport at curve carriers;
\item \textbf{V5} the $5$-letter type calculus at $g=2$: the level
profile of all Schur constituents of $\otimes^5V$; only
$\wedge^4V\otimes V$ at level 1;
\item \textbf{V6} the channel-typing shadow of Lemma 1 (the shift+twist
type arithmetic on random graded models);
\item \textbf{V7} the CM shadow: for the Gaussian letters
($D=\Q(i)$-module structure), the $D$-invariant middle-typed tensors of
$\otimes^3$ computed exactly; the non-CM control fires;
\item \textbf{V8} the NS-boundary at self-products:
$[\wedge^3V\cap(\NS\otimes V)]=\theta\wedge V$ by exact RREF at
$g=2,3$;
\item \textbf{V9} hostile controls: the sabotage battery --- broken
transport (a non-correspondence map), wrong twist ($\delta$ of one-off
codimension), fake volume (a non-top exterior power), fake CM (a random
involution) --- every control must FIRE with the predicted failure.
\end{itemize}
All loops bounded before launch (enumerations $\le2^{12}$, matrices
$\le12\times12$, instances $\le400$, runtime $<60$\,s, zero floats,
\texttt{Fraction}/int only), the vacuity scan (ledger-107 patterns)
clean, deterministic double-run byte-identical.

\section{The honest scope}

Nothing in this round discharges (R2$'$), (R3) or (R4) at general
input. What changed is the \emph{trajectory}, for the second consecutive
round: the residual family (R1) that looked like an unbounded
general-input frontier is gone --- replaced by three named objects with
strictly smaller total content. The results that did the work are
unconditional (S7, S8a on its stated stratum, S8b, S9i--ii) or
classical-sector identifications (S9iii). The very-general closure of
S8a(c) is a \emph{stratum} result, stated with its stratum condition on
the face (the MT condition is checkable, and the stratum is
Zariski-dense); it is not a claim on the special strata, where the
merged Weil sector lives and remains open. Per the campaign's banned
substitutes discipline (PROMPT \S3): curve-product strata results are
reported as \emph{sector/stratum} closures with the residual named ---
never as the conjecture. The falsifiable line, sharpened: (R2$'$) one
beyond-divisor balanced class on one $E\times A$ at a Weil-type input
(the merged sector); (R3) the $K$-condition at one general input beyond
the channels; (R4) the fold-birth at one transcendental input with the
input-sensitivity gate firing.

\end{document}
